\section{Architecture}
\label{sec:arch}

The policy of Section~\ref{sec:method} is defined over causal
information alone, but a software implementation on top of a
block-level cache cannot reconstruct the precise event timing that the
scoring contract requires, as the vLLM study in
Section~\ref{sec:sw-vllm} will confirm. A hardware realization
addresses three structural limitations that software cannot.
First, the controller must observe every request arrival, gap onset,
and gap closure with cycle-level fidelity, which a near-memory
placement provides without interrupt or polling overhead.
Second, eviction and migration decisions must complete within the
tool-gap window so that the promoted KV is already in the fast tier
when the next request arrives, which demands a deterministic pipeline
rather than a software thread.
Third, eviction and tiering are two reads of one ranking, and a single
register file eliminates the staleness that any dual-engine design
would introduce.

This section maps the algorithmic policy to a synthesizable
control-plane IP. Section~\ref{sec:place} defines the controller
scope and system interface. Section~\ref{sec:org} derives the unified
control plane from the interleaving structure of agent workloads.
Section~\ref{sec:quant} locks the integer contract through a format
sweep. Section~\ref{sec:micro} maps the result to a pipelined event
core. Section~\ref{sec:dse} explores the session-capacity design
space and selects the design point. Fig.~\ref{fig:arch} shows the
resulting architecture.

\begin{figure*}[t]
\centering
\includegraphics[width=\textwidth]{fig03_architecture.pdf}
\caption{\unison hardware architecture. The scheduling core integrates
\spear ranking and \tide migration into a unified pipeline beside the
two-tier KV hierarchy.}
\label{fig:arch}
\end{figure*}

\subsection{Controller Scope and System Placement}
\label{sec:place}

\kcmu is a session-level control-plane IP that ranks residency in a
shared KV pool. It does not compute attention, store KV payload, or
manage the HBM PHY. Runtime software writes events, namely request
arrivals, gap onsets, gap closures, and completions, into a CSR event
FIFO, and the controller returns a victim index and logical DMA
descriptors. Physical address mapping, KV arrays, the HBM PHY, and
the on-chip network remain external. The IP therefore
owns~\eqref{eq:scoreq} and~\eqref{eq:budget}, not the bytes that
those decisions move.

Near-memory placement beside the hierarchy is motivated by the event
interface. A software scheduler that polls block-cache timestamps
approximates gap onset with millisecond granularity at best, and
Section~\ref{sec:sw-vllm} shows that this approximation inverts the
gap signal when tool waits approach the decode window. A controller
wired to the memory hierarchy observes the same events at cycle
granularity without software intervention, which is the structural
advantage that Section~\ref{sec:hw-place} quantifies under placement
variation.

\subsection{Unified Control Plane}
\label{sec:org}

Agent workloads interleave request and gap events at fine granularity
rather than in long homogeneous phases. A dual design that pairs a
\spear engine with a \tide engine under periodic synchronization is
the textbook decomposition, but it is structurally unsound under this
interleaving pattern. Eviction and promotion are two reads of one
ranking, not two independent computations, so a second engine is
either stale or redundant.

Let $R_{\mathrm{u}}$ denote the unified register file and
$R_{\mathrm{d}}$ the pair of private views in a dual design. The dual
stores
\begin{equation}
  \lvert R_{\mathrm{d}}\rvert=2\lvert R_{\mathrm{u}}\rvert
  \label{eq:rf}
\end{equation}
and applies a stale snapshot at every decision between synchronization
points. If the last sync was at event $e_{k-m}$, the dual victim is
\begin{equation}
  v_{\mathrm{d}}
  =\arg\max_{s\in\mathcal{L}}
    \mathrm{score}_{q}\!\bigl(s;\,R_{\mathrm{d}}(e_{k-m})\bigr),
  \label{eq:stale}
\end{equation}
whereas the unified victim uses $R_{\mathrm{u}}(e_{k})$.
Tab.~\ref{tab:org} confirms the structural prediction. Even with
every-event synchronization the dual design diverges on SRAM
placement, and with reduced synchronization the stale view causes
hit-rate drops of up to 5.4\,pp and inconsistent AMAT signs across
traces, confirming that a unified register file is not merely
preferable but structurally necessary.

% Auto-generated by paper/tables/gen_tab5_org.py
\begin{table}[t]
\centering
\renewcommand{\arraystretch}{1.12}
\setlength{\tabcolsep}{3pt}
\scriptsize
\begin{threeparttable}
\caption{Architectural ablation of a unified control plane versus independent dual~IPs under varying synchronization granularity.}
\label{tab:org}
\begin{tabular*}{\columnwidth}{@{\extracolsep{\fill}}lrrrrrr@{}}
\toprule
\textbf{Trace} & \textbf{Vict$\neq$}\tnote{a} & \textbf{Prom$\neq$} &
\textbf{$\Delta$HR}\tnote{b} & \textbf{$\Delta$SRAM} &
\textbf{$\Delta$Pre} & \textbf{$\Delta$AMAT}\tnote{c} \\
 & (\%) & (\%) & (\%) & (\%) & (\%) & (\%) \\
\midrule
\rowcolor{gray!10}
\multicolumn{7}{c}{\textbf{Dual IPs copy $R$ after every request and gap}} \\
SWE/Qw3 & 0 & 0 & $+0.0$ & $+0.0$ & $+0.0$ & $+0.0$ \\
SWE/Dev & 0 & 0 & $+0.0$ & $+0.0$ & $+0.0$ & $+0.0$ \\
SWE/Ge4 & 0 & 0 & $+0.1$ & $+0.0$ & $-0.0$ & $+0.0$ \\
GAIA/Qw3 & 0 & 0 & $+0.0$ & $+0.5$ & $+0.8$ & $-1.7$ \\
GAIA/Dev & 0 & 0 & $+0.0$ & $-0.2$ & $+0.8$ & $-4.7$ \\
GAIA/Ge4 & 0 & 0 & $+0.2$ & $-1.3$ & $+1.0$ & $-1.8$ \\
\rowcolor{gray!10}
\multicolumn{7}{c}{\textbf{Dual IPs copy $R$ every 8 requests}} \\
SWE/Qw3 & 64 & 55 & $-2.2$ & $-0.4$ & $-1.2$ & $+1.2$ \\
SWE/Dev & 68 & 62 & $-1.9$ & $+0.0$ & $-1.1$ & $+1.1$ \\
SWE/Ge4 & 53 & 43 & $-5.4$ & $+0.0$ & $-1.9$ & $+2.0$ \\
GAIA/Qw3 & 35 & 50 & $-3.5$ & $-0.7$ & $-3.2$ & $+6.4$ \\
GAIA/Dev & 18 & 68 & $-0.2$ & $+1.0$ & $+0.9$ & $-5.0$ \\
GAIA/Ge4 & 39 & 45 & $-4.5$ & $-1.2$ & $-4.1$ & $+7.3$ \\
\rowcolor{gray!10}
\multicolumn{7}{c}{\textbf{Dual IPs copy $R$ only at gap start}} \\
SWE/Qw3 & 3 & 29 & $-0.0$ & $-0.1$ & $-0.0$ & $+0.0$ \\
SWE/Dev & 2 & 43 & $+0.2$ & $+0.2$ & $+0.0$ & $-0.0$ \\
SWE/Ge4 & 4 & 16 & $+0.4$ & $+0.0$ & $+0.1$ & $-0.1$ \\
GAIA/Qw3 & 5 & 39 & $+0.9$ & $-0.2$ & $+1.2$ & $-2.5$ \\
GAIA/Dev & 13 & 77 & $-3.4$ & $+0.7$ & $-1.7$ & $+9.0$ \\
GAIA/Ge4 & 4 & 29 & $+0.8$ & $-1.9$ & $+1.2$ & $-2.1$ \\
\bottomrule
\end{tabular*}
\begin{tablenotes}\scriptsize
\item[a] Vict$\neq$ and Prom$\neq$ are the shares of copied evictions and promotions that pick a different session.
\item[b] $\Delta$HR, $\Delta$SRAM, and $\Delta$Pre are percentage-point shifts versus the live table.
\item[c] $\Delta$AMAT is the percent AMAT change, and either sign represents a consistency failure.
\end{tablenotes}
\end{threeparttable}
\end{table}


The unified pipeline therefore keeps one live register file and
sequences \esp before \okt in the same datapath, paying
$\lvert R_{\mathrm{u}}\rvert$ rather than~\eqref{eq:rf}.

\subsection{Fixed-Point Quantization of the Scoring Contract}
\label{sec:quant}

Realizing the floating-point policy in synthesizable logic requires an
integer image of~\eqref{eq:esp} whose ranking fidelity is
indistinguishable from the reference. The local complement
$\sigma(t)$ rather than the cumulative product-limit survival is the
key enabler, because a single ROM lookup avoids the sequential
dependence chain that a cumulative estimator would impose on the
pipeline. Let $b$ be a fractional width and $\Lambda_b=2^{b}-1$.
The unsigned quantizer
\begin{equation}
  Q_{b}(x)
  =\mathrm{clip}\!\left(
      \bigl\lfloor x\,\Lambda_b+\tfrac12\bigr\rfloor,\;
      0,\;\Lambda_b
    \right)
  \label{eq:uq}
\end{equation}
maps $x\in[0,1]$ onto $\mathrm{UQ}0.b$. Hazard and survival then
become
\begin{align}
  h_{q}
  &=Q_{b_{h}}\!\bigl(h(\min(t,50))\bigr),
  \label{eq:hq}
  \\
  \sigma_{q}
  &=\max\!\bigl(
      \bigl\lfloor\varepsilon\Lambda_{b_{h}}+\tfrac12\bigr\rfloor,\;
      \Lambda_{b_{h}}-h_{q}
    \bigr).
  \label{eq:survq}
\end{align}
A time divisor $\tau\in\{1,10^{3},10^{6}\}$ stores the gap as
\begin{equation}
  g_{q}
  =\mathrm{clip}\!\left(
      \bigl\lfloor g_s/\tau\bigr\rfloor,\;
      0,\;2^{b_{g}}-1
    \right),
  \label{eq:gapq}
\end{equation}
so $\tau=1$ is nanoseconds, $\tau=10^{3}$ is microseconds, and
$\tau=10^{6}$ is milliseconds. The two score terms are integers,
\begin{align}
  T_{1}
  &=\left\lfloor
      \frac{g_{q}\,\Lambda_{b_{h}}}{\sigma_{q}}
    \right\rfloor,
  \label{eq:t1}
  \\
  T_{2}
  &=\left\lfloor
      \frac{(\Lambda_{b_{h}}-\sigma_{q})\,P_{\mathbb{Z}}}
           {\Lambda_{b_{h}}}
    \right\rfloor,
  \label{eq:t2}
\end{align}
and the hardware score is the saturating sum
\begin{equation}
  \mathrm{score}_{q}(s)
  =\mathrm{sat}_{b_{s}}\!\bigl(T_{1}+T_{2}\bigr),
  \qquad
  b_{s}=64,
  \label{eq:scoreq}
\end{equation}
with $P_{\mathbb{Z}}=P/\tau$ so that the multiplicative penalty
remains an affine image of $(1-\sigma)P$ after the time scaling. The
EMA in~\eqref{eq:ema} is the same fraction in integer arithmetic,
\begin{equation}
  g_{q}
  \leftarrow
  \left\lfloor
    \frac{3\,g_{q}^{\mathrm{new}}+7\,g_{q}}{10}
  \right\rfloor.
  \label{eq:ema-int}
\end{equation}
A completed session writes $2^{b_{s}}-1$.

The contract is selected on a pre-registered grid of 180 format
points that vary $b_{h}\in\{8,10,12\}$,
$b_{g}\in\{24,32,64\}$, $\tau\in\{1,10^{3},10^{6}\}$, divider type,
and penalty form. Ranking fidelity is measured on 7265 real eviction
snapshots from the six traces and three capacity envelopes of
Section~\ref{sec:sweval}. The hard gates are top-1 victim agreement
${\ge}0.99$ and maximum hit-rate drift ${\le}0.20$\,pp. The sweep
selects
\begin{equation}
  (b_{h},b_{g},b_{s},\tau,P_{\mathbb{Z}})
  =(12,32,64,10^{3},5\times10^{8}),
  \label{eq:lock}
\end{equation}
at which the aggregate Kendall $\tau$ exceeds 0.998 and the worst-cell
hit-rate drift is 0.192\,pp. Twelve-bit hazard is the minimum width that passes the ranking gate,
microsecond storage balances register width against fidelity, and an
exact divider is preferred over a reciprocal LUT because the LUT ties
fidelity while adding a table of depth $\Lambda_{b_{h}}$. Section~\ref{sec:dse} shows that
the bits ranking demanded are essentially free relative to the session
register file.

\subsection{Event Core Microarchitecture}
\label{sec:micro}

The unified pipeline must evaluate~\eqref{eq:scoreq} for every live
session within the tool-gap window so that the eviction or migration
decision is ready before the next request arrives. A parallel
comparator tree would close in $O(1)$ cycles but scale as $O(N)$ in
area. A pipelined linear scan trades latency for area, closing in
$N{+}5$ cycles at a cost that grows with the register file rather
than with dedicated comparator logic. The five-stage depth maps
directly to the arithmetic of the quantized score, with the
critical-path divider of~\eqref{eq:t1} dictating the pipeline depth.

The stages are
\begin{equation}
  \begin{aligned}
    \mathrm{S0}&:\;
      (t,g_{q},\mathrm{done})\leftarrow R[s],\\
    \mathrm{S1}&:\;
      h_{q}\leftarrow\mathrm{ROM}[t],\;
      \sigma_{q}\leftarrow~\eqref{eq:survq},\\
    \mathrm{S2}&:\;
      T_{1}\leftarrow~\eqref{eq:t1},\\
    \mathrm{S3}&:\;
      T_{2}\leftarrow~\eqref{eq:t2},\\
    \mathrm{S4}&:\;
      \mathrm{score}_{q}\leftarrow~\eqref{eq:scoreq},\;
      \text{update running }\arg\max.
  \end{aligned}
  \label{eq:pipe}
\end{equation}
A full scan of $N$ live sessions takes
\begin{equation}
  L_{\mathrm{scan}}=N+5
  \label{eq:scan}
\end{equation}
cycles after pipeline fill, yielding a decision latency of
\begin{equation}
  \tau(N)=\frac{N+5}{f_{\mathrm{clk}}}.
  \label{eq:tau}
\end{equation}
The \tide budget uses the same microsecond gap,
\begin{equation}
  \mathcal{B}_{q}
  =\bigl\lfloor\Delta/\tau\bigr\rfloor\cdot B_{\mu},
  \label{eq:budgetq}
\end{equation}
where $B_{\mu}$ is tokens per microsecond.
Cosimulation of the locked format against the floating-point reference
is bit-exact on all 1\,937 events and 845 decisions of the GAIA/Gemma4
Edge-Tight trace.

Given the pipeline and the quantization contract, the remaining
architectural degree of freedom is the session capacity~$N$.

\subsection{Session-Capacity Design-Space Exploration}
\label{sec:dse}

Session capacity~$N$ governs both decision latency
through~\eqref{eq:tau} and register-file area. To evaluate the policy
side of this tradeoff, we define Agent Serving Effectiveness (ASE) as
the six-dataset geometric mean
\begin{equation}
  \mathrm{ASE}(N)=
  \Biggl(\prod_{d=1}^{6}
    \frac{\mathrm{HR}_d}{\mathrm{HR}_d^{\max}}
    \cdot
    \frac{\mathrm{AMAT}_d^{\min}}{\mathrm{AMAT}_d}
    \cdot
    \frac{\mathrm{TTFT}_d^{\min}}{\mathrm{TTFT}_d}
  \Biggr)^{1/18}
  \label{eq:ase}
\end{equation}
inside the capacity sweep, where the inner product runs over three
normalized metrics and the outer over six traces.

Fig.~\ref{fig:dse} reports the exploration on the Yosys/Nangate45
flow used to rank variants by relative cell count. ASE saturates near
$N{=}64$ where trace concurrency plateaus, while decision latency
from~\eqref{eq:tau} remains well below the median tool gap at that
point. The capacity knee is therefore $N{=}64$. Absolute ASE values,
latency measurements, and the Design Compiler PPA at this point are
reported in Section~\ref{sec:hweval}.

\begin{figure}[t]
\centering
\begin{subfigure}{0.49\linewidth}
\includegraphics[width=\linewidth]{fig08a_pareto_2d.pdf}
\caption{Cell count and ASE versus $N_{\mathrm{sess}}$}
\label{fig:dse-2d}
\end{subfigure}\hfill
\begin{subfigure}{0.49\linewidth}
\includegraphics[width=\linewidth]{fig08b_n64_zoom.pdf}
\caption{Structural variants at $N{=}64$}
\label{fig:dse-n64}
\end{subfigure}
\caption{Design-space exploration of session capacity and structural
variants on the Yosys/Nangate45 synthesis flow.}
\label{fig:dse}
\end{figure}

Fig.~\ref{fig:dse-n64} details the $N{=}64$ column. The five live
configurations occupy a narrow cell-count band, and the selected
design point retains the exact divider of~\eqref{eq:t1} for bit-exact
scoring against the software reference, because the reciprocal-LUT
alternative ties ranking fidelity without improving the end-to-end
gate.

Sections~\ref{sec:sweval} and~\ref{sec:hweval} evaluate the policy
and report synthesis results.
